\chapter{From Kitaev Chain to Topological Quantum Computation}

\section{Kitaev Chain}

A toy model proposed by Kitaev (2001) that exhibits unpaired Majorana zero modes at its ends. It captures two essential physical ingredients:
\begin{itemize}
    \item The $U(1)$ symmetry $\hat{c}_j \mapsto \ee^{\ii \phi} \hat{c}_j$, corresponding to electric charge conservation, must be broken down to a $\Z_2$ fermion parity symmetry $\oc_j \mapsto -\oc_j$. This is achieved by proximity-coupling the wire to an $s$-wave superconductor.
    \item A triplet (p-wave) pairing potential between electrons of the same spin is required. A spinless (or spin-polarized) wire adjacent to a 3D $s$-wave superconductor can effectively realize such $p$-wave pairing once spin-orbit coupling is accounted for.
\end{itemize}

Consider a 1D chain of $N$ sites, each of which can be either empty or occupied by a spinless electron. The Hamiltonian is
\begin{align}
    \oH = \sum_{j=1}^{N} \left[ \underbrace{-t\!\left(\oc^{\dag}_{j}\,\oc_{j+1} + \oc^{\dag}_{j+1}\,\oc_{j}\right)}_{\text{hopping}} \underbrace{-\, \mu \!\left(\oc^{\dag}_{j}\,\oc_{j} - \tfrac{1}{2}\right)}_{\text{chemical potential}} + \underbrace{\Delta\,\oc_{j}\,\oc_{j+1} + \Delta^{*}\oc^{\dag}_{j+1}\,\oc^{\dag}_{j}}_{\text{p-wave pairing}} \right],
    \label{eq:kitaev_H}
\end{align}
where $t > 0$ is the hopping amplitude, $\mu$ is the chemical potential (half-filling is $\mu = 0$), and $\Delta = |\Delta|\ee^{\ii\phi} \in \mathbb{C}$ is the superconducting order parameter inherited from the substrate. The operators obey the canonical anti-commutation relations
\begin{align}
    \anticomm{\oc_j}{\oc_k^{\dag}} = \delta_{jk}, \qquad \anticomm{\oc_j}{\oc_k} = 0.
\end{align}
By a global gauge transformation $\oc_j \mapsto \ee^{\ii\phi/2}\,\oc_j$, the phase of $\Delta$ can be absorbed, so we set $\Delta \in \R_{>0}$ hereafter without loss of generality. Open boundary conditions (OBC) are used in Secs.~\ref{sec:zeromodes}--\ref{sec:finitesize}; periodic boundary conditions (PBC) are used in Sec.~\ref{sec:fourier}.

%---------------------------------------------------------------------------
\subsection{Majorana Operators}
%---------------------------------------------------------------------------

\begin{inter}
Real Majorana particles can propagate freely through space; the quasi-Majorana modes in the Kitaev chain are different---they are bound to the ends of the wire and cannot move. Their topological protection comes precisely from being spatially separated.
\end{inter}

Every complex fermionic mode $(\oc_j, \oc_j^{\dag})$ can be decomposed into two \emph{Majorana} operators $\gamma_{2j-1}$ and $\gamma_{2j}$ via
\begin{align}
    \oc_j = \frac{\gamma_{2j-1} + \ii\,\gamma_{2j}}{2}, \qquad
    \oc_j^{\dag} = \frac{\gamma_{2j-1} - \ii\,\gamma_{2j}}{2},
    \label{eq:maj_def}
\end{align}
which satisfy the Majorana algebra
\begin{align}
    \anticomm{\gamma_a}{\gamma_b} = 2\delta_{ab}, \qquad \gamma_a^{\dag} = \gamma_a.
    \label{eq:maj_algebra}
\end{align}
In particular $\gamma_a^2 = 1$: a Majorana operator squares to the identity. The \emph{inverse} of \eqref{eq:maj_def} expresses the Majorana operators back in terms of the original fermions:
\begin{align}
    \gamma_{2j-1} = \oc_j + \oc_j^{\dag}, \qquad
    \gamma_{2j} = -\ii\!\left(\oc_j - \oc_j^{\dag}\right) = \ii\!\left(\oc_j^{\dag} - \oc_j\right).
    \label{eq:maj_inv}
\end{align}
One may verify \eqref{eq:maj_algebra} directly from the CAR: e.g.\ $\{\gamma_{2j-1},\gamma_{2j-1}\} = 2(\oc_j + \oc_j^{\dag})^2 = 2\{\oc_j,\oc_j^{\dag}\} = 2$.

A useful identity relates the local occupation number to a Majorana bilinear:
\begin{align}
    \oc_j^{\dag}\,\oc_j = \frac{1 + \ii\,\gamma_{2j-1}\,\gamma_{2j}}{2},
    \label{eq:occ_maj}
\end{align}
so the chemical potential term becomes $-\mu(\oc_j^{\dag}\oc_j - \tfrac{1}{2}) = -\tfrac{\ii\mu}{2}\,\gamma_{2j-1}\,\gamma_{2j}$.

\paragraph{Hamiltonian in the Majorana language.}
Substituting \eqref{eq:maj_def} into \eqref{eq:kitaev_H} and using the Majorana algebra, one can evaluate each term. Writing $a_j \equiv \gamma_{2j-1}$ and $b_j \equiv \gamma_{2j}$ for brevity:

\noindent\textit{Hopping:}
\begin{align}
    -t\!\left(\oc_j^{\dag}\oc_{j+1} + \oc_{j+1}^{\dag}\oc_j\right)
    = \frac{\ii t}{2}\!\left(\gamma_{2j}\,\gamma_{2j+1} - \gamma_{2j-1}\,\gamma_{2j+2}\right).
\end{align}

\noindent\textit{Pairing (real $\Delta$):}
\begin{align}
    \Delta\!\left(\oc_j\,\oc_{j+1} + \oc_{j+1}^{\dag}\oc_j^{\dag}\right)
    = \frac{\ii\Delta}{2}\!\left(\gamma_{2j}\,\gamma_{2j+1} + \gamma_{2j-1}\,\gamma_{2j+2}\right).
\end{align}

Combining all three pieces, the full Hamiltonian reads
\begin{align}
    \boxed{
    \oH = \sum_{j=1}^{N}\left[
    -\frac{\ii\mu}{2}\,\gamma_{2j-1}\,\gamma_{2j}
    +\frac{\ii(\Delta+t)}{2}\,\gamma_{2j}\,\gamma_{2j+1}
    +\frac{\ii(\Delta-t)}{2}\,\gamma_{2j-1}\,\gamma_{2j+2}
    \right].}
    \label{eq:H_maj}
\end{align}

\paragraph{Two instructive limits.}
\eqref{eq:H_maj} makes the topology transparent by considering two extreme parameter choices:

\begin{description}
\item[Trivial phase] $t = \Delta = 0$, $\mu \neq 0$:
\begin{align}
    \oH_{\rm triv} = -\frac{\ii\mu}{2}\sum_{j=1}^{N}\gamma_{2j-1}\,\gamma_{2j}.
\end{align}
The Majorana operators pair \emph{within} each physical site (Fig.~\ref{fig:majorana_chain}a). Both Majoranas at every site form a single ordinary fermion; there are no leftover, unpaired modes.

\item[Topological phase] $\mu = 0$, $\Delta = t > 0$:
\begin{align}
    \oH_{\rm top} = \ii t\sum_{j=1}^{N-1}\gamma_{2j}\,\gamma_{2j+1}.
    \label{eq:H_top_special}
\end{align}
The Majorana operators pair \emph{between} adjacent sites (Fig.~\ref{fig:majorana_chain}b). The operators $\gamma_1$ (leftmost) and $\gamma_{2N}$ (rightmost) are completely absent from $\oH_{\rm top}$. They commute with the Hamiltonian, cost zero energy, and constitute the two \emph{Majorana zero modes}.
\end{description}

%---------------------------------------------------------------------------
\subsection{Fourier Transformation to K-space}
\label{sec:fourier}
%---------------------------------------------------------------------------

We now work with periodic boundary conditions ($\oc_{N+1} \equiv \oc_1$) and pass to momentum space. Define the Fourier modes
\begin{align}
    \oc_j = \frac{1}{\sqrt{N}}\sum_{k}\ee^{\ii kj}\,\oc_k, \qquad
    k = \frac{2\pi n}{N},\quad n = 0,1,\ldots,N-1.
\end{align}
Expressing \eqref{eq:kitaev_H} in terms of $\oc_k$:
\begin{align}
    \oH = \sum_k\left[\xi_k\,\oc_k^{\dag}\oc_k
    + \frac{1}{2}\!\left(d_k\,\oc_k\,\oc_{-k}
    + d_k^{*}\,\oc_{-k}^{\dag}\oc_k^{\dag}\right)\right] + \text{const},
\end{align}
where the normal-state dispersion and pairing amplitude are
\begin{align}
    \xi_k = -2t\cos k - \mu, \qquad d_k = -2\ii\Delta\sin k.
    \label{eq:xi_d}
\end{align}

\paragraph{Bogoliubov--de Gennes (BdG) formalism.}
Introduce the Nambu (particle-hole) spinor
\begin{align}
    \Psi_k = \begin{pmatrix}\oc_k \\ \oc_{-k}^{\dag}\end{pmatrix},
\end{align}
which doubles the degrees of freedom in exchange for a simple matrix structure. The Hamiltonian takes the standard BdG form
\begin{align}
    \oH = \frac{1}{2}\sum_k\Psi_k^{\dag}\,\mathcal{H}_{\rm BdG}(k)\,\Psi_k + \text{const},
    \label{eq:H_BdG}
\end{align}
with the $2\times 2$ BdG matrix
\begin{align}
    \mathcal{H}_{\rm BdG}(k)
    = \begin{pmatrix} \xi_k & d_k \\ d_k^{*} & -\xi_k \end{pmatrix}
    = \begin{pmatrix} -2t\cos k - \mu & -2\ii\Delta\sin k \\ 2\ii\Delta\sin k & 2t\cos k+\mu \end{pmatrix}.
    \label{eq:BdG_matrix}
\end{align}
In terms of Pauli matrices $\boldsymbol{\sigma} = (\sigma_x,\sigma_y,\sigma_z)$:
\begin{align}
    \mathcal{H}_{\rm BdG}(k) = \underbrace{(-2t\cos k - \mu)}_{d_z(k)}\,\sigma_z
    + \underbrace{2\Delta\sin k}_{d_y(k)}\,\sigma_y
    = \mathbf{d}(k)\cdot\boldsymbol{\sigma},
    \label{eq:BdG_dvec}
\end{align}
where $\mathbf{d}(k) = (0,\,d_y(k),\,d_z(k))$ is a 2D vector confined to the $yz$-plane.

\paragraph{Quasiparticle spectrum.}
The BdG eigenvalues come in $\pm$ pairs (particle-hole symmetry):
\begin{align}
    E_{\pm}(k) = \pm\,|\mathbf{d}(k)| = \pm\sqrt{(2t\cos k + \mu)^{2} + 4\Delta^{2}\sin^{2}k}.
    \label{eq:spectrum}
\end{align}
The gap closes at $E_\pm(k)=0$ precisely when $\mathbf{d}(k) = \mathbf{0}$:
\begin{itemize}
    \item At $k = 0$: gap closes when $\mu = -2t$.
    \item At $k = \pm\pi$: gap closes when $\mu = +2t$.
\end{itemize}
These two conditions define the \textbf{phase boundaries} $|\mu| = 2t$, across which the topology changes (see Sec.~\ref{sec:topinv}).

\paragraph{Particle-hole symmetry.}
The BdG Hamiltonian satisfies
\begin{align}
    \sigma_x\,\mathcal{H}_{\rm BdG}^{*}(-k)\,\sigma_x = -\mathcal{H}_{\rm BdG}(k),
\end{align}
reflecting the redundancy built into the Nambu doubling: for every Bogoliubov quasiparticle at $(k,E)$ there is a partner at $(-k,-E)$.

%---------------------------------------------------------------------------
\subsection{Topological Invariants}
\label{sec:topinv}
%---------------------------------------------------------------------------

\paragraph{Symmetry class.}
The Kitaev chain has particle-hole symmetry $\mathcal{C}$ (with $\mathcal{C}^2 = +1$) and broken time-reversal symmetry, placing it in \emph{symmetry class D}. According to the periodic table of topological insulators and superconductors, 1D class-D systems are classified by a $\Z_2$ invariant.

\paragraph{Winding number.}
A concrete $\mathbb{Z}$ winding number can be extracted from the trajectory of $\mathbf{d}(k)$ in the $yz$-plane as $k$ sweeps the Brillouin zone $[-\pi,\pi]$. Define the unit vector $\hat{\mathbf{d}}(k) = \mathbf{d}(k)/|\mathbf{d}(k)|$ and the angle $\theta(k) = \arctan[d_y(k)/d_z(k)]$. The winding number is
\begin{align}
    \nu = \frac{1}{2\pi}\int_{-\pi}^{\pi}\dd k\,\frac{\dd\theta}{\dd k}
    = \frac{1}{2\pi}\oint_{\rm BZ} \frac{d_z\,\dd d_y - d_y\,\dd d_z}{d_z^2+d_y^2} \in \mathbb{Z}.
    \label{eq:winding}
\end{align}
Geometrically, $\nu$ counts how many times $\mathbf{d}(k)$ encircles the origin. The $\Z_2$ topological invariant is $(-1)^{\nu}$, but since $\nu$ can only be $0$ or $\pm 1$ for the Kitaev chain, the integer already captures the physics.

\paragraph{Phase diagram.}
One can evaluate \eqref{eq:winding} analytically by tracking the endpoints $\mathbf{d}(0) = (-2t - \mu,\, 0)$ and $\mathbf{d}(\pm\pi) = (2t - \mu,\, 0)$:
\begin{itemize}
    \item If $\mu = 0$: $\mathbf{d}(0) = (-2t, 0)$ points in $-z$, while $\mathbf{d}(\pm\pi) = (2t, 0)$ points in $+z$. As $k$ increases from $0$ to $\pi$, the $y$-component $d_y = 2\Delta\sin k > 0$ lifts $\mathbf{d}$ above the axis and it winds once around the origin: $\nu = 1$.
    \item If $|\mu| > 2t$: both endpoints $\mathbf{d}(0)$ and $\mathbf{d}(\pm\pi)$ lie on the \emph{same} side of the origin ($d_z$ never changes sign), so the path winds zero times: $\nu = 0$.
\end{itemize}
The general result is
\begin{align}
    \nu = \begin{cases} 1 & |\mu| < 2t \quad \text{(topological phase)},\\ 0 & |\mu| > 2t \quad \text{(trivial phase)},\end{cases}
\end{align}
with phase transitions at $|\mu| = 2t$ where the bulk gap \eqref{eq:spectrum} closes.

\begin{theorem}[Bulk--boundary correspondence]
When $|\mu| < 2t$ ($\nu = 1$), a semi-infinite chain possesses exactly one Majorana zero mode exponentially localized at its boundary. For OBC, a chain of $N$ sites has two such modes, one at each end, that are exponentially split by $\delta E \sim \ee^{-N/\xi}$ where $\xi$ is the bulk localization length.
\end{theorem}

\paragraph{Pfaffian invariant.}
Alternatively, using Kitaev's original formulation, the $\Z_2$ invariant can be extracted from the Pfaffian of the antisymmetric BdG matrix at the time-reversal-invariant momenta $k^* = 0, \pi$:
\begin{align}
    (-1)^{\nu} = \operatorname{sgn}\!\left[\operatorname{Pf}\!\left(\mathcal{H}_{\rm BdG}(0)\right)\right]\cdot\operatorname{sgn}\!\left[\operatorname{Pf}\!\left(\mathcal{H}_{\rm BdG}(\pi)\right)\right].
\end{align}
At $k^* = 0$: $\operatorname{Pf}(\mathcal{H}(0)) = d_z(0) = -2t - \mu$. At $k^* = \pi$: $\operatorname{Pf}(\mathcal{H}(\pi)) = d_z(\pi) = 2t - \mu$. Therefore
\begin{align}
    (-1)^{\nu} = \operatorname{sgn}\!\left[(2t + \mu)(2t - \mu)\right] = \operatorname{sgn}\!\left[4t^2 - \mu^2\right],
\end{align}
recovering the same phase boundary $|\mu| = 2t$.

%---------------------------------------------------------------------------
\subsection{Zero Modes}
\label{sec:zeromodes}
%---------------------------------------------------------------------------

We now work with OBC and focus on the Majorana zero modes that emerge in the topological phase. These are the central object of interest for topological quantum computation.

\paragraph{Exact solution at the special point.}
At $\mu = 0$, $\Delta = t$, the Hamiltonian \eqref{eq:H_top_special} contains only the inter-site couplings $\gamma_{2j}\gamma_{2j+1}$. The operators
\begin{align}
    \Gamma_L \equiv \gamma_1, \qquad \Gamma_R \equiv \gamma_{2N}
    \label{eq:zm_special}
\end{align}
do not appear in $\oH_{\rm top}$ and therefore commute with it: $[\oH_{\rm top},\Gamma_{L/R}] = 0$. They are \emph{exact} zero modes with vanishing energy cost. Note that $\Gamma_L$ is a single-site operator localized at the left end ($j=1$), and $\Gamma_R$ at the right end ($j=N$) — perfectly localized, with zero correlation length.

\paragraph{Non-local fermionic qubit.}
The two Majorana zero modes combine into a single complex fermion
\begin{align}
    \hat{f} = \frac{\Gamma_L + \ii\,\Gamma_R}{2}, \qquad
    \hat{f}^{\dag} = \frac{\Gamma_L - \ii\,\Gamma_R}{2},
    \label{eq:nonlocal_fermion}
\end{align}
satisfying $\{\hat{f},\hat{f}^{\dag}\} = 1$. Since $[\oH,\hat{f}] = 0$, the occupation number $\hat{n}_f = \hat{f}^{\dag}\hat{f}$ is conserved and the ground-state manifold is \emph{two-fold degenerate}:
\begin{align}
    |GS_0\rangle \quad (n_f = 0),\qquad |GS_1\rangle = \hat{f}^{\dag}|GS_0\rangle \quad (n_f = 1).
\end{align}
These two states have \emph{opposite} fermion parity and the same energy. Crucially, the quantum information encoded in the basis $\{|GS_0\rangle, |GS_1\rangle\}$ is stored \textbf{non-locally}: any local operator acting only near one end cannot distinguish the two states, making the qubit immune to local perturbations.

\paragraph{Generic parameter case.}
Away from the special point but still in the topological phase ($|\mu| < 2t$, $\Delta \neq 0$), the zero modes are no longer perfectly localized at single sites; instead they decay exponentially into the bulk. The BdG zero-mode equation $\mathcal{H}_{\rm BdG}\Phi_0 = 0$ with the ansatz $\Phi_0(j) \propto \lambda^j$ leads to a characteristic polynomial whose roots $\lambda_{1,2}$ satisfy $|\lambda_{1,2}| < 1$. The Majorana wavefunction takes the form
\begin{align}
    \phi_{L}(j) \propto \lambda_1^{\,j} - \lambda_2^{\,j}, \qquad
    |\phi_L(j)| \sim \ee^{-j/\xi},
\end{align}
with localization length
\begin{align}
    \xi^{-1} = -\ln|\lambda_{\min}|,
\end{align}
where $|\lambda_{\min}| < 1$ is the smaller root magnitude. The localization length diverges at the phase boundaries $|\mu| \to 2t$ and $\Delta \to 0$, as the bulk gap closes.

\begin{definition}[Majorana zero mode]
A \emph{Majorana zero mode} (MZM) is a self-adjoint operator $\Gamma = \Gamma^{\dag}$ that (i) commutes with the Hamiltonian $[\oH,\Gamma] = 0$, (ii) has zero eigenvalue (i.e.\ costs no energy), and (iii) is exponentially localized near a boundary or defect.
\end{definition}

\paragraph{Stability.}
The MZMs are protected as long as:
\begin{enumerate}[label=(\roman*)]
    \item The bulk gap remains open (no quantum phase transition);
    \item Fermion parity is conserved (no single-electron tunneling from the environment);
    \item The two ends are far enough apart that their wavefunctions do not overlap appreciably.
\end{enumerate}
Any local perturbation that preserves these conditions cannot split the ground-state degeneracy.

%---------------------------------------------------------------------------
\subsection{Finite Size Chain}
\label{sec:finitesize}
%---------------------------------------------------------------------------

For a chain of finite length $N$, the left and right Majorana wavefunctions have exponential tails that penetrate into the bulk and eventually overlap near the center of the chain. This overlap lifts the ground-state degeneracy.

\paragraph{Energy splitting.}
Let $\Phi_L(j)$ and $\Phi_R(j)$ be the BdG wavefunctions of the left and right MZMs respectively. Their overlap integral is
\begin{align}
    M = \sum_j \Phi_L^{\dag}(j)\,\mathcal{H}_{\rm BdG}\,\Phi_R(j) \approx A\,\ee^{-N/\xi}\cos\!\left(k_F N + \phi_0\right),
    \label{eq:splitting}
\end{align}
where $k_F$ is an effective Fermi wavevector associated with the oscillatory part of the decaying wavefunction, and $\phi_0$ is a phase. The resulting energy splitting between the two formerly degenerate ground states is $\delta E = 2|M|$.

\paragraph{Oscillatory behavior.}
The cosine factor in \eqref{eq:splitting} has a striking consequence: as $N$ increases, the sign of $\delta E$ can alternate, meaning the ground state can switch between even and odd fermion parity with period $\sim \pi/k_F$ in system size. This oscillation is a direct probe of the Fermi surface geometry and is observable in transport or spectroscopy measurements.

\paragraph{Special point.}
At $\mu = 0$, $\Delta = t$, the MZMs are perfectly localized ($\xi = 0$ formally, or more precisely one site), so their wavefunctions have exactly zero overlap for any $N \geq 2$:
\begin{align}
    \delta E\big|_{\mu=0,\,\Delta=t} = 0 \quad \text{for all } N.
\end{align}
This is confirmed by exact diagonalization: the lowest two eigenvalues of the BdG matrix are numerically zero at machine precision.

\paragraph{Scaling near the phase boundary.}
Approaching the phase boundary from the topological side, the bulk gap $\Delta_{\rm gap} \sim ||\mu| - 2t|$ closes linearly, and the localization length diverges as
\begin{align}
    \xi \sim \frac{v_F}{\Delta_{\rm gap}} \sim \frac{1}{||\mu| - 2t|},
\end{align}
where $v_F = 2t\sin k_F$ is an effective Fermi velocity. As $\xi \to \infty$, finite-size effects \eqref{eq:splitting} become significant for any fixed chain length, and the MZMs hybridize strongly.

\paragraph{Numerical determination of the topological phase.}
A practical diagnostic is to compute the lowest positive eigenvalue $E_{\min}(N)$ of the BdG matrix for OBC as a function of $\mu$:
\begin{itemize}
    \item In the topological phase: $E_{\min}(N) \to 0$ exponentially fast as $N \to \infty$.
    \item In the trivial phase: $E_{\min}(N)$ converges to a finite gap value.
\end{itemize}
The value of $\mu$ at which $E_{\min}(N)$ collapses to zero identifies the phase boundary, which sharpens with increasing $N$.
